\documentclass{article}
\usepackage{amsmath,amssymb}
\usepackage{xurl}
\usepackage[hidelinks]{hyperref}
% Shared commands for notes in docs/paper-gaps/.

\DeclareUnicodeCharacter{2080}{\ifmmode{}_{0}\else\textsubscript{0}\fi}
\DeclareUnicodeCharacter{2081}{\ifmmode{}_{1}\else\textsubscript{1}\fi}
\DeclareUnicodeCharacter{2082}{\ifmmode{}_{2}\else\textsubscript{2}\fi}
\DeclareUnicodeCharacter{2099}{\ifmmode{}_{n}\else\textsubscript{n}\fi}

\newcommand{\Lean}{\textsc{Lean}}
\ExplSyntaxOn
\NewDocumentCommand{\leanid}{m}{
  \ifmmode
    \textsf{\detokenize{#1}}
  \else
    \group_begin:
    \urlstyle{sf}
    \tl_set:Nn \l_tmpa_tl {#1}
    \tl_replace_all:Nnn \l_tmpa_tl {\_} {_}
    \exp_args:NV \path \l_tmpa_tl
    \group_end:
  \fi
}
\ExplSyntaxOff
\providecommand{\tr}{\operatorname{tr}}
\newcommand{\tnleanIssueBase}{https://github.com/LionSR/TNLean/issues/}
\newcommand{\tnleanPrBase}{https://github.com/LionSR/TNLean/pull/}
\newcommand{\tnleanDocBase}{https://sirui-lu.com/TNLean/docs/}
\newcommand{\ghissue}[1]{\href{\tnleanIssueBase#1}{GitHub issue \##1}}
\newcommand{\ghpr}[1]{\href{\tnleanPrBase#1}{PR \##1}}
\newcommand{\doclink}[1]{\href{\tnleanDocBase#1.html}{\path{#1}}}

% Dirac notation and the identity operator, as in blueprint/src/macros/common.tex.
% \providecommand keeps note-local definitions authoritative.
\usepackage{dsfont}
\providecommand{\ket}[1]{|#1\rangle}
\providecommand{\bra}[1]{\langle #1|}
\providecommand{\braket}[2]{\langle #1 | #2 \rangle}
\providecommand{\ketbra}[2]{|#1\rangle\!\langle #2|}
\providecommand{\Id}{\mathds{1}}
\providecommand{\one}{\mathds{1}}
\providecommand{\C}{\mathbb{C}}
\providecommand{\MN}[1]{M_{#1}(\C)}
\providecommand{\MD}{M_D(\C)}

% Verdict marker: \gapnote{<kind>}{<status>}. Kind is the policy
% classification (clarification, local-correction, scope-restriction,
% unfaithful, false-source, open-gap); status is open, wip (elimination
% underway), resolved, or historical. Machine-read by the site index;
% prints a verdict line.
\newcommand{\gapnote}[2]{\par\noindent\textbf{Verdict:} #1 (#2).\par}

\title{Identity-Ancilla Transport of MPU Symmetries}
\author{}
\date{2026-08-29}
\begin{document}
\maketitle
\gapnote{scope-restriction}{open}

\section{Source definition}

Cirac--P\'erez-Garc\'ia--Schuch--Verstraete define two MPU tensors to be
equivalent under a symmetry when identity ancillas and a common blocking make
them strictly equivalent under that symmetry
\cite[lines~1356--1366]{Cirac2017MPU}. The definition does not state how an
action on the original physical dimension determines an action after adjoining
an identity ancilla.

\section{Missing transport}

Let $d>0$ be the original physical dimension, let $N>0$ be the chain length,
and let $x>1$ be an ancilla dimension. The canonical sitewise rearrangement is
\begin{align}
    R_{d,x,N}\colon
    (\mathbb C^d)^{\otimes N}\otimes(\mathbb C^x)^{\otimes N}
     & \longrightarrow (\mathbb C^{dx})^{\otimes N}.
    \label{eq:mpu_symmetry_ancilla_transport_sitewise_rearrangement}
\end{align}
It induces the operator embedding
\begin{align}
    \iota_x(X)
     & = R_{d,x,N}(X\otimes I_{x^N})R_{d,x,N}^{-1}.
    \label{eq:mpu_symmetry_ancilla_transport_operator_embedding}
\end{align}
Under these nondegeneracy assumptions, the map in
(\ref{eq:mpu_symmetry_ancilla_transport_operator_embedding}) is not onto the
enlarged operator algebra. Consequently, an action at dimension $d$ can
constrain the enlarged action on $\iota_x(X)$ but cannot determine its values
on arbitrary operators at dimension $dx$. A strict-equivalence path can leave
the image of $\iota_x$, so an action defined only on that image is insufficient.

The formalized restricted construction therefore takes actions at every
physical dimension as supplied data and assumes
\begin{align}
    \iota_x(\mathcal S_{d,N}(X))
     & = \mathcal S_{dx,N}(\iota_x(X)).
    \label{eq:mpu_symmetry_ancilla_transport_semiconjugacy}
\end{align}
The symmetry family supplies the condition in
(\ref{eq:mpu_symmetry_ancilla_transport_semiconjugacy}), while the stabilized
equivalence also fixes one ambient virtual bond dimension along the path.\footnote{Formal declarations:
    \leanid{MPOTensor.FiniteChainOperatorSymmetryFamily} supplies the
    dimension-indexed actions and semiconjugacy law;
    \leanid{MPOTensor.EquivalentUnderSymmetry} defines the fixed-bond
    coherent-family equivalence.}

\section{Additional common ancillas}

The preliminary equivalence convention chooses positive coprime ancilla
sizes $p_a,p_b$ with $p_a d_a=p_b d_b$
\cite[lines~706--724]{Cirac2017MPU}. Later, the source compares two tensors
of the same physical dimension by adjoining an identity ancilla of the same
dimension $d$ to each tensor
\cite[lines~2210--2215]{Cirac2017MPU}. For $d>1$, this common pair is not
coprime. In fact, when $d_a=d_b>0$, dimension balance forces $p_a=p_b$,
and coprimality then forces $p_a=p_b=1$.

Thus the current coprime-ancilla equivalence predicates cannot express this
later common-ancilla construction directly. This is a further restriction,
independent of the missing symmetry extension in
(\ref{eq:mpu_symmetry_ancilla_transport_semiconjugacy}). A source-level
stabilized equivalence must allow a common additional ancilla after the
minimal coprime enlargement, or equivalently permit arbitrary positive
balanced ancilla sizes. Its symmetry actions must be defined on the full
enlarged operator algebra, and its paths must address the virtual-dimension
restriction separately. The fixed-dimension one-sided swap comparison does
not discharge any of these stabilization requirements.

\section{Verdict}

\textbf{Verdict: scope restriction.} The coherent-family definition is a
well-defined strengthened variant, not the unrestricted source definition.
The source-level definition remains unformalized until the intended actions on
the full enlarged operator algebras are specified without adding the
semiconjugacy hypothesis to the cited statement, and the common-ancilla
construction is admitted without the minimal coprime-pair restriction.

\bibliographystyle{alpha}
\bibliography{references}
\end{document}